
\documentclass{article}
\usepackage{colortbl}
\usepackage{makecell}
\usepackage{multirow}
\usepackage{supertabular}

\begin{document}

\newcounter{utterance}

\centering \large Interaction Transcript for game `referencegame', experiment `random\_grids', episode 17 with reka{-}flash{-}3.1{-}t0.0.
\vspace{24pt}

{ \footnotesize  \setcounter{utterance}{1}
\setlength{\tabcolsep}{0pt}
\begin{supertabular}{c@{$\;$}|p{.15\linewidth}@{}p{.15\linewidth}p{.15\linewidth}p{.15\linewidth}p{.15\linewidth}p{.15\linewidth}}
    \# & $\;$A & \multicolumn{4}{c}{Game Master} & $\;\:$B\\
    \hline

    \theutterance \stepcounter{utterance}  
    & & \multicolumn{4}{p{0.6\linewidth}}{
        \cellcolor[rgb]{0.9,0.9,0.9}{
            \makecell[{{p{\linewidth}}}]{
                \texttt{\tiny{[P1$\langle$GM]}}
                \texttt{You are given three grids, where each of them is 5 by 5 in size.} \\
\texttt{Grids have empty cells marked with "▢" and filled cells marked with "X".} \\
\texttt{Your task is to generate a referring expression that best describes the target grid while distinguishing it from the two other distractor grids.} \\
\texttt{The first grid is the target grid, and the following two grids are the distractors.} \\
\\ 
\texttt{Target grid:} \\
\\ 
\texttt{X X □ X □} \\
\texttt{□ X X □ X} \\
\texttt{X □ □ X □} \\
\texttt{□ X □ □ X} \\
\texttt{X □ X □ X} \\
\\ 
\texttt{Distractor grid 1:} \\
\\ 
\texttt{X X □ □ X} \\
\texttt{□ X X □ X} \\
\texttt{X □ X X □} \\
\texttt{□ X □ X X} \\
\texttt{X □ □ X X} \\
\\ 
\texttt{Distractor grid 2:} \\
\\ 
\texttt{X □ X X □} \\
\texttt{□ X □ X X} \\
\texttt{X X □ □ X} \\
\texttt{□ X X □ X} \\
\texttt{X □ X □ □} \\
\\ 
\texttt{Instruction: Describe the target grid.} \\
\texttt{Generate the referring expression starting with the tag "Expression: " for the given target grid. Omit any other text.} \\
            }
        }
    }
    & \\ \\

    \theutterance \stepcounter{utterance}  
    & \multicolumn{4}{p{0.6\linewidth}}{
        \cellcolor[rgb]{0.9,0.9,0.9}{
            \makecell[{{p{\linewidth}}}]{
                \texttt{\tiny{[P1$\rangle$GM]}}
                \texttt{<reasoning>} \\
\texttt{Let's tackle this problem. I need to generate a referring expression that describes the target grid and distinguishes it from the two distractor grids. The target grid is the first one, and the other two are distractors. The grids are 5x5, with Xs and □s.} \\
\\ 
\texttt{First, I should look at the target grid and see what unique patterns or features it has that the distractors don't. Let me write down the target grid again to visualize it:} \\
\\ 
\texttt{Row 1: X X □ X □} \\
\texttt{Row 2: □ X X □ X} \\
\texttt{Row 3: X □ □ X □} \\
\texttt{Row 4: □ X □ □ X} \\
\texttt{Row 5: X □ X □ X} \\
\\ 
\texttt{Now, let's check Distractor 1:} \\
\\ 
\texttt{Row 1: X X □ □ X} \\
\texttt{Row 2: □ X X □ X} \\
\texttt{Row 3: X □ X X □} \\
\texttt{Row 4: □ X □ X X} \\
\texttt{Row 5: X □ □ X X} \\
\\ 
\texttt{Distractor 2:} \\
\\ 
\texttt{Row 1: X □ X X □} \\
\texttt{Row 2: □ X □ X X} \\
\texttt{Row 3: X X □ □ X} \\
\texttt{Row 4: □ X X □ X} \\
\texttt{Row 5: X □ X □ □} \\
\\ 
\texttt{Hmm. Let's compare each row and column for differences. Maybe look for specific positions where the target differs from the others. Let's go row by row.} \\
\\ 
\texttt{Starting with Row 1 of Target: X X □ X □. Distractor 1's Row 1 is X X □ □ X. So in Target, positions 3 is □, 4 is X, 5 is □. In Distractor 1, positions 3 is □, 4 is □, 5 is X. So the difference here is at position 4 and 5. Distractor 2's Row 1 is X □ X X □. So positions 2 is □, 3 is X, 4 is X, 5 is □. So Target's Row 1 has two Xs at start, then □, X, □. Distractor 1 has two X} \\
            }
        }
    }
    & & \\ \\

    \theutterance \stepcounter{utterance}  
    & & & \multicolumn{2}{p{0.3\linewidth}}{
        \cellcolor[rgb]{0.9,0.9,0.9}{
            \makecell[{{p{\linewidth}}}]{
                \texttt{\tiny{[GM$|$GM]}}
                \texttt{Invalid generated expression} \\
            }
        }
    }
    & & \\ \\

\end{supertabular}
}

\end{document}
